\chapter{METODE DAN RANCANGAN PENELITIAN}

\section{Deskripsi Umum dan Status Rancangan}

Penelitian ini merancang evaluasi empiris terkontrol atas pengelolaan imperatif terotomasi berbasis \texttt{kubectl} (Lingkungan A, disebut manual dalam judul) dan GitOps berbasis ArgoCD (Lingkungan B). Beban formal adalah InvenioRDM nyata beserta dependensi multi-service, bukan aplikasi NGINX pengganti. Yang dibandingkan adalah dua prosedur operasional pada konfigurasi tertentu, bukan keunggulan universal suatu paradigma.

Revisi metodologis ini merupakan draf untuk tinjauan pemilik penelitian dan pembimbing. Eksperimen formal belum disetujui dan eksekusi langsung instrumen dinonaktifkan. Kelulusan uji offline tidak membuktikan kesiapan klaster, audit, SMTP, kapasitas, ataupun pemulihan data aplikasi. Proposal yang diterima sebelumnya dan artefak historis tetap dapat ditelusuri melalui riwayat versi; perubahan definisi tidak mengubah makna hasil historis menjadi hasil formal.

\section{Metodologi dan Kesetaraan Perbandingan}

Rancangan mengikuti evaluasi empiris terkontrol dengan pengamatan kondisi dasar, satu intervensi, dan pengamatan pascaintervensi \citep{wohlin2012experimentation}. Kedua lingkungan dibangun secara sengaja pada satu klaster; rancangan bukan penugasan acak klaster dari populasi \citep{cook1979quasi}. Unit analisis adalah satu injeksi yang diterima pada satu lingkungan dalam satu pengulangan.

Kedua lingkungan menggunakan sumber deployment, image digest, inventaris, kontrak bidang konfigurasi, kebijakan HPA, dan aturan pengukuran yang setara. Lingkungan manual berada di luar kepemilikan ArgoCD. Urutan lingkungan dan skenario diacak dalam blok; hanya satu injeksi diperbolehkan setelah reset sebelumnya terverifikasi. Tidak ada eksekusi semua mutasi tanpa reset di antaranya.

Pemeriksa imperatif membandingkan keadaan aktual dengan acuan yang dibekukan, menyimpan bukti perbedaan, lalu menjalankan koreksi tetap. Tidak ada penundaan untuk reaksi atau pengetikan manusia. Identitas skenario tidak diberikan kepada mekanisme koreksi. Peneliti membakukan prosedur dan meninjau bukti, bukan menjadi subjek pengukuran kinerja manusia; NASA-TLX dan waktu identifikasi akar masalah manusia tidak menjadi metrik rancangan ini. Prosedur pemulihan tetap mengikuti prinsip runbook yang dapat diperiksa \citep{schlette2024playbook}.

\begin{figure}[H]
  \centering
  \begin{tikzpicture}[
    box/.style={draw, rounded corners, align=center, text width=3.5cm, minimum height=10mm, font=\small},
    flow/.style={->, thick}, node distance=7mm and 8mm]
    \node[box] (source) {Sumber InvenioRDM\\revisi dan render dibekukan};
    \node[box, below left=of source] (manual) {A: otomasi kubectl\\HPA normal};
    \node[box, below right=of source] (gitops) {B: ArgoCD\\HPA normal};
    \node[box, below=24mm of source, text width=5cm] (measure) {Evaluator bersama\\R1 bidang dan inventaris\\R2 kejadian teramati\\R3 bukti audit};
    \draw[flow] (source) -- (manual);
    \draw[flow] (source) -- (gitops);
    \draw[flow] (manual) -- (measure);
    \draw[flow] (gitops) -- (measure);
  \end{tikzpicture}
  \caption{Rancangan apparatus dengan HPA aktif pada kedua lingkungan.}
  \label{fig:arsitektur-eksperimen}
\end{figure}

Gambar~\ref{fig:arsitektur-eksperimen} membedakan mekanisme deployment dari evaluator. HPA merupakan pengendali bersama, sehingga responsnya tidak diatribusikan kepada ArgoCD.

\begin{figure}[H]
  \centering
  \begin{tikzpicture}[
    box/.style={draw, rounded corners, align=center, text width=5cm, minimum height=8mm, font=\small},
    flow/.style={->, thick}, node distance=5mm]
    \node[box] (review) {Tinjauan metode dan validasi offline};
    \node[box, below=of review] (ready) {Gerbang kapasitas, audit, aplikasi\\dan pembakuan sumber};
    \node[box, below=of ready] (baseline) {Verifikasi baseline dan aktifkan kolektor};
    \node[box, below=of baseline] (inject) {Satu injeksi dan bukti penerimaan};
    \node[box, below=of inject] (observe) {Amati sampai kejadian atau batas waktu};
    \node[box, below=of observe] (reset) {Reset lengkap terverifikasi\\sebelum pengulangan berikutnya};
    \node[box, below=of reset] (analysis) {Analisis per skenario\\A dan F terpisah dari efikasi ArgoCD};
    \draw[flow] (review) -- (ready);
    \draw[flow] (ready) -- (baseline);
    \draw[flow] (baseline) -- (inject);
    \draw[flow] (inject) -- (observe);
    \draw[flow] (observe) -- (reset);
    \draw[flow] (reset) -- (analysis);
  \end{tikzpicture}
  \caption{Bagan alir rancangan penelitian; pelaksanaan formal tetap menunggu persetujuan.}
  \label{fig:bagan-alir-eksperimen}
\end{figure}

Gambar~\ref{fig:bagan-alir-eksperimen} menempatkan gerbang kesiapan sebelum injeksi. Kegagalan dan data tersensor disimpan, bukan dihapus agar hanya pemulihan sukses yang dilaporkan.

\section{Lingkungan dan Provenans Deployment}

Konteks deployment yang dilaporkan pemilik adalah tiga node, 24 inti CPU, sekitar 23 Gi memori, Kubernetes 1.32.6 dengan RKE2/containerd, dan Ubuntu 24.04. Nilai ini menggantikan target lama Kubernetes 1.29+ dan Ubuntu 22.04; bukan hasil pemeriksaan live baru dalam revisi ini. Sebelum pengumpulan formal, versi runtime, alokasi node, dan revisi yang benar-benar diterapkan harus direkam ulang.

Sumber otoritatif adalah repositori \url{https://github.com/vityasyyy/invenio-rdm-gitops}. Revisi referensi yang diperiksa ialah \nolinkurl{6cacea6642a7c1690243d3b54130fa08216a9595}; penunjuknya disimpan pada \path{experiments/deployment-source.json}. Application \path{argocd/apps/invenio-bootstrap.yaml} memilih \path{k8s/apps/invenio/kustomization.yaml}, yang merujuk sumber daya dan image digest. Revisi referensi bukan pernyataan bahwa klaster sedang menjalankan revisi itu. Tidak dibuat salinan manifest Invenio editable di repositori penelitian.

Sebelum eksperimen, catat revisi sumber aktual, hasil render beserta SHA-256, image digest, versi ArgoCD dan chart, inventaris objek, penyesuaian namespace/isolasi yang disetujui, serta kepemilikan pengendali. Hasil render merupakan artefak bukti tak berubah, bukan sumber konfigurasi kedua. Dependensi bersama dibedakan dari sumber daya yang diisolasi agar dua lingkungan tidak saling mengubah data.

Konfigurasi sumber \path{k8s/infra/argocd/patches/argocd-cm.yaml} menetapkan \texttt{timeout.reconciliation} sebesar 30 detik, berbeda dari target proposal lama 60 detik. Interval rekonsiliasi repositori ini tidak sama dengan watch kejadian Kubernetes, retry self-heal, atau interval evaluator. Seluruh interval dan perilaku aktual harus dicatat serta dibekukan sebelum formal, bukan disesuaikan setelah melihat hasil.

\section{Intervensi dan Kepemilikan Pengendali}

Pemilihan kategori mempertimbangkan defect konfigurasi Kubernetes dan praktik operasional \citep{zhang2026configuration,pohjola2025mitigating}. Tabel~\ref{tab:matriks-eksperimen} membedakan kontrol HPA, drift terkelola, dan batas visibilitas.

\begin{table}[H]
  \centering
  \small
  \caption{Matriks skenario dan batas interpretasi}
  \label{tab:matriks-eksperimen}
  \begin{tabularx}{\textwidth}{|c|>{\raggedright\arraybackslash}p{3.3cm}|X|}
    \hline
    Skenario & Intervensi & Kepemilikan dan analisis \\
    \hline
    A & Perubahan replika Deployment & HPA aktif pada kedua lengan; kontrol respons HPA, bukan efikasi ArgoCD \\
    \hline
    B & Perubahan image & Bidang Deployment terpilih; image digest acuan dibekukan \\
    \hline
    C & Modifikasi ConfigMap & Kesetaraan objek dan adopsi konfigurasi proses dibedakan \\
    \hline
    D & Penghapusan sumber daya & Verifikasi penciptaan ulang dan seluruh bidang; bukan sekadar rollout sehat \\
    \hline
    E & Patch requests/limits & Seluruh bidang terpilih dibandingkan dengan aturan identik \\
    \hline
    F & Objek pada namespace salah & Batas visibilitas objek tak terkelola; tidak diasumsikan dipangkas ArgoCD \\
    \hline
    G & Perubahan Secret sintetis & Isolasi dari pengendali SealedSecret wajib ditinjau; nilai tidak diterbitkan \\
    \hline
  \end{tabularx}
\end{table}

Pada A, HPA web dan worker tetap memiliki replika pada kedua lingkungan. Sumber yang diperiksa mengabaikan \texttt{/spec/replicas} pada perbandingan ArgoCD. Karena itu, tidak munculnya \textit{OutOfSync} bukan waktu deteksi nol, dan kembalinya replika bukan bukti self-heal ArgoCD. Nilai acuan dinamis berasal dari keputusan/status HPA yang teramati, bukan angka tetap tiga replika. Kebijakan, metrik, rentang, serta jendela stabilisasi HPA dan penghormatan kepemilikan saat sync harus diverifikasi. Hasil A dilaporkan sebagai kontrol HPA, tidak digabungkan ke estimasi efikasi B--E dan G.

F dilaporkan tersendiri karena objek di luar cakupan Application tidak otomatis menjadi sasaran prune/self-heal. Objek uji harus dibangun dari spesifikasi tersanitasi, bukan menyalin UID, resourceVersion, managedFields atau status objek server ke namespace lain. Penerimaan injeksi dan pembersihan karantina wajib dibuktikan.

Pada G, Secret produksi yang dihasilkan SealedSecret tidak boleh diperlakukan sebagai koreksi ArgoCD murni. Kontrak Secret sintetis dan pengendalinya harus disetujui sebelum formal. Perubahan laten ImageUpdater juga harus dibekukan dalam apparatus yang disetujui. Tidak ada izin mengubah pengendali produksi melalui rancangan ini.

\section{Variabel dan Definisi Operasional}

\subsection{R1: Kesetaraan Konfigurasi}

Kontrak inventaris dan bidang terpilih dibekukan dari sumber yang sama untuk kedua lingkungan. Sumber daya cocok hanya jika ada dan seluruh bidang terpilihnya sama dengan acuan. Kontrak mencakup image, requests/limits, isi ConfigMap, kesetaraan Secret, serta bidang Service dan objek lain yang relevan.

\[
R1 = \frac{\text{jumlah sumber daya dengan seluruh bidang terpilih cocok}}{\text{jumlah sumber daya dalam inventaris acuan}} \times 100\%.
\]

Status Healthy/Synced bukan pengganti kesetaraan. Kesehatan dan adopsi runtime dilaporkan terpisah. Bidang dinamis HPA dikecualikan dari kesetaraan statis secara eksplisit, bukan dihitung sebagai drift terhadap angka replika tetap. Objek ekstra/karantina dicatat terpisah dari penyebut. Objek hilang pada pembacaan lengkap adalah ketidakcocokan; pembacaan API gagal atau bidang tidak diperiksa adalah kegagalan bukti, bukan skor nol yang sah.

R1 dicatat pada baseline, pengamatan pertama pascainjeksi, dan pascakoreksi. Jika drift sudah pulih sebelum pembacaan, kondisi sesaatnya ditandai tidak teramati. Secret dibandingkan di memori; hanya identitas bidang dan boolean kesetaraan disimpan, tanpa nilai atau hash yang dapat digunakan untuk menebak isinya.

\subsection{R2: Deteksi, Bukti Deviasi, dan Pemulihan}

Semua durasi dihitung dari satu jam monotonik kolektor dalam detik dengan pecahan; UTC disimpan untuk audit. Kolektor harus aktif sebelum injeksi. Interval nominal evaluator satu detik merupakan target validasi, bukan kemampuan live yang telah dibuktikan. Interval aktual dan jeda pembacaan dicatat, terpisah dari jadwal mekanisme koreksi.

\begin{itemize}
  \item $t_0$: pengiriman permintaan injeksi; $t_a$: penerimaan respons API. Injeksi sah hanya jika API berhasil dan perubahan bidang diverifikasi dari respons atau pembacaan balik.
  \item $t_d$: sinyal pertama mekanisme pengelolaan yang relevan berdasarkan perbedaan konfigurasi. Pada B--E/G, sumbernya pemeriksa imperatif atau status perbedaan ArgoCD sesuai lingkungan. Rollout gagal bukan definisi drift; rollout sehat dapat tetap memiliki image, konfigurasi, atau limits yang salah.
  \item $t_e$: waktu pertama bukti tersimpan mengidentifikasi objek dan bidang yang berbeda. Diff menunjukkan deviasi, bukan penyebab kausal; status Healthy dan selesainya perintah diagnosis tidak menghasilkan waktu akar masalah.
  \item $t_r$: pengamatan pertama seluruh inventaris dan bidang kembali cocok, dikonfirmasi pada pembacaan berikutnya. Pada A, konfirmasi harus mencakup keadaan dinamis HPA; pada F, ketiadaan objek karantina diperiksa tersendiri.
\end{itemize}

Durasi utama adalah $t_d-t_0$, $t_e-t_0$, dan $t_r-t_0$. Pada A, sinyal HPA dan respons replika diberi label pengendali HPA, bukan deteksi ArgoCD. Tidak adanya sinyal sampai batas pengamatan dicatat tersensor kanan dengan waktu hilang, tidak diisi nol atau disamakan dengan $t_0$. Pemulihan konfigurasi tidak membuktikan layanan sehat atau data aplikasi pulih.

\subsection{R3: Bukti Keterlacakan}

R3 dideskripsikan dengan rubrik yang sama untuk pelaku, objek/bidang perubahan, waktu/urutan, dan tindakan/hasil koreksi. Setiap unsur diberi status ada, sebagian, atau tidak ada disertai referensi bukti tersanitasi. Git mencatat perubahan deklarasi, bukan otomatis pelaku mutasi \texttt{kubectl} di luar Git. Karena itu, audit API dibutuhkan pada kedua lengan untuk merekonstruksi mutasi tersebut; kelengkapan tidak diasumsikan karena suatu alat digunakan. R3 tidak diuji sebagai perbedaan yang hasilnya telah ditentukan dan tidak menghasilkan klaim akar masalah.

\section{Gerbang Baseline dan Reset}

Sebelum injeksi, inventaris sumber daya, pod dan node wajib tidak kosong serta lengkap; semua pod layanan harus Running dan Ready. CPU dan memori setiap node harus di bawah 70\%. Memori didefinisikan sebagai working set metrics-server dibagi allocatable node, bukan metrik host yang dapat dipertukarkan. Metrik tidak tersedia, field tidak terbaca, inventaris parsial, atau status pengendali tidak diketahui menggagalkan gerbang secara tertutup. Kelulusan sampel tunggal tidak menggantikan bukti jendela beban stabil.

Kesiapan juga mencakup bukti audit, SMTP yang benar-benar dapat digunakan, isolasi, sinkronisasi jam, kepemilikan HPA, beban latar, serta pemulihan aplikasi/data/queue yang telah ditinjau. Audit 22 Oktober 2026 menemukan node worker02 melewati 70\%; konteks kesiapan yang dilaporkan pemilik tetap sekitar 71\%. Kapasitas dua lingkungan Invenio nyata belum terbukti. Audit API belum ditetapkan, SMTP masih placeholder, dan private-draft seed serta recovery aplikasi belum tervalidasi. Tidak ada prasyarat tersebut yang dianggap selesai oleh perubahan instrumen.

Bootstrap, migrasi, fixtures, indexing dan deklarasi queue harus dipisahkan dari repair rutin. Sync perbaikan yang memicu setup dapat menambah kerja dan mengubah baseline. Antarmuka lifecycle wajib mengembalikan exit code, waktu, identitas tindakan dan postcondition; kegagalan tak terduga tidak boleh ditelan. ConfigMap/Secret yang dipakai melalui envFrom memerlukan bukti adopsi runtime terpisah, bukan hanya kesetaraan objek.

Setelah satu pengulangan, reset wajib membuktikan seluruh inventaris/bidang cocok, tidak ada objek karantina, pod siap, kondisi sumber daya memenuhi batas, keadaan HPA sesuai, serta data, queue dan konfigurasi proses kembali ke acuan. Rollout sehat saja tidak cukup. Kegagalan satu tindakan koreksi diteruskan ke pemanggil dan pengulangan berikutnya dilarang sampai reset lengkap terverifikasi. Tidak ada pengosongan queue, modifikasi data, restore atau migrasi langsung yang diizinkan oleh naskah ini.

\section{Instrumen dan Protokol Pengukuran}

Empat entrypoint \texttt{baseline-check.sh}, \texttt{drift-inject.sh}, \texttt{metrics-collect.sh}, dan \texttt{recovery-runbook.sh} kini memvalidasi bundel bukti tersimpan secara offline. Pemeriksa bersama menolak bukti hilang/tidak konsisten, menggunakan aturan R1 identik, membedakan sinyal pengendali dari rollout, serta mempertahankan waktu tersensor. Uji fixture mencakup sukses, kegagalan dan batas; fixture bukan hasil pengukuran aplikasi. Adapter polling, injeksi, pemulihan/reset langsung, ekspor render, integrasi lifecycle, audit dan SMTP belum selesai; pengujian formal tidak boleh dimulai sebelum adapter tersebut ditinjau dan divalidasi.

Urutan formal yang direncanakan ialah: (1) bekukan metode, sumber, batas waktu dan aturan analisis; (2) simpan baseline serta aktifkan kolektor; (3) kirim satu injeksi dan verifikasi penerimaan; (4) catat sinyal pengendali, bukti deviasi, kesetaraan dan kesehatan sampai pulih atau batas pengamatan; (5) simpan kegagalan dan hasil tersensor; (6) verifikasi reset lengkap sebelum pengulangan berikutnya. Tidak ada eksekusi formal dalam revisi ini.

\section{Pengulangan dan Analisis Statistik}

Target deskriptif tetap $n=5$ per skenario per lingkungan untuk E1/R1 dan $n=10$ untuk E2/R2; E3/R3 deskriptif. A adalah kontrol HPA dan F batas visibilitas, keduanya dilaporkan terpisah dari B--E/G. Menjumlahkan skenario yang heterogen tidak membuktikan daya uji per skenario. Target ini bukan hasil analisis power yang tervalidasi dan tidak memberi jaminan kecukupan sampel \citep{cohen1988statistical}.

Laporkan jumlah injeksi sah, kejadian teramati/tersensor, kegagalan instrumen/pemulihan/reset, median, rentang dan distribusi per skenario. Waktu tidak teramati tetap kosong dengan status sensor dan batas pengamatan. Kegagalan API/injeksi/infrastruktur dibedakan dari tidak terjadinya deteksi; data tidak dibuang tanpa jejak. Rencana uji tetap Mann--Whitney U dua sisi pada $\alpha=0{,}05$, dengan ukuran efek Cohen's $d$ atau Cliff's $\delta$ sebagaimana proposal sebelumnya. Penerapannya memerlukan observasi lengkap serta asumsi yang sesuai; revisi ini tidak menetapkan pilihan ukuran efek atau keluarga uji baru. Jika ada penyensoran, analisis yang mengabaikannya tidak dianggap perbandingan lengkap; keputusan analisis data tersensor dan penyelarasan rincian proposal--skripsi memerlukan tinjauan pemilik/pembimbing sebelum formal.

\section{Ancaman terhadap Validitas}

Validitas internal dipengaruhi perbedaan jadwal pengendali, tekanan node, I/O latar, dependensi bersama, hooks lifecycle, dan kepemilikan HPA/SealedSecret/ImageUpdater. Kendali dan pencatatannya membatasi bias, tidak menjamin satu-satunya perbedaan adalah paradigma abstrak. Validitas konstruk dibatasi kontrak bidang, resolusi pengamatan, bukti audit, serta perbedaan kesetaraan konfigurasi dengan kesehatan dan pemulihan aplikasi. A dan F tidak dapat digunakan untuk klaim efikasi self-heal ArgoCD.

Validitas eksternal dibatasi satu klaster, InvenioRDM, ArgoCD dan konfigurasi yang akhirnya digunakan. Surrogat S1/S2 serta hasil mock historis bukan sampel formal dan tidak memperluas inferensi. Reliabilitas memerlukan provenans, jam kolektor, bukti lengkap, dan reset yang dapat diulang. Sampai tinjauan metode serta seluruh gerbang dipenuhi, belum tersedia dasar untuk kesimpulan komparatif empiris.
